\documentclass[11pt,a4paper]{article}
\usepackage{amsmath,amssymb,amsthm,mathtools,bm}
\usepackage[utf8]{inputenc}
\usepackage[T1]{fontenc}
\usepackage{amsfonts}
\usepackage{graphicx}
\usepackage{xcolor}
\usepackage{textcomp}
\usepackage[margin=1in]{geometry}
\usepackage{lmodern}
\usepackage{microtype}

\usepackage{booktabs,array,longtable}
\usepackage[hidelinks]{hyperref}
\usepackage{physics}
\usepackage{setspace}
\setstretch{1.06}
\usepackage{tikz}
\usetikzlibrary{arrows.meta,positioning,calc,fit,decorations.pathmorphing}
% Tables and Figures
\usepackage{float}
\usepackage{import}
\usepackage{tabularx}
\usepackage{enumitem}
\usepackage{appendix}

%-------------------------- SST minimal macro prelude (compile-safe) --------------------
% Vector swirl velocity + clock/vorticity symbols (aligned with Rosetta)
\newcommand{\omegas}{\boldsymbol{\omega}_{\!\mkern-2mu\scriptstyle\boldsymbol{\circlearrowleft}}}
\newcommand{\rc}{r_c}
\newcommand{\rhoF}{\rho_{\!f}}
\newcommand{\rhoE}{\rho_{\!E}}
\newcommand{\rhoM}{\rho_{\!m}}
\newcommand{\rhocore}{\rho_{\text{core}}}
\newcommand{\boxedstatement}[1]{\fbox{\parbox{0.86\linewidth}{\centering #1}}}
%----------------------------------------------------------------------------------------

\title{Toward a Canonical Atomic Torsion-Bridge in Swirl-String Theory\\
\large v0.9 Circulation-Foundational Atomic Torsion Bridge}
\author{Omar Iskandarani \\
Independent Researcher, Groningen, The Netherlands \\
\texttt{info@omariskandarani.com} \\
ORCID: 0009-0006-1686-3961}
\date{March 2026}

\begin{document}
\maketitle

\begin{abstract}
We present a self-contained hydrodynamic derivation of atomic structure within Swirl-String Theory (SST). By modeling the vacuum as an incompressible, inviscid condensate, we replace abstract quantum postulates with deterministic fluid-topological mechanics. The primitive circulation quantum $\Gamma_0$ acts as the fundamental driver of the atomic sector. Matter corresponds to closed, topologically stable vortex loops (trefoils) governed by a macroscopic continuous envelope, while electromagnetic radiation is modeled as a transverse torsion field. We derive the hydrogenic Bohr radius as a thermodynamic equilibrium surface balancing dynamic swirl pressure and vacuum tension. Expanding to multi-electron systems, we demonstrate that the Pauli exclusion principle and the $N_n = 2n^2$ Aufbau shell capacity emerge directly from the fluid's cavitation limit and the discrete delay-induced phase locking of the macroscopic circulation field. A photoionization benchmark is formulated via a branch-resolved continuum transition rate, yielding an operator-level handedness asymmetry for the hydrogenic source state. The framework systematically recovers standard atomic phenomenology while predicting falsifiable symmetry-breaking behaviors in chiral photoionization.

\vspace{0.5cm}
\noindent \textbf{Keywords:} Swirl-String Theory, Topological Fluid Dynamics, Aufbau Principle, Delay-Differential Dynamics, Fermionic Exchange
\end{abstract}

\section{Aim and canonical constraints}
The objective is not to replace the successful mathematics of atomic physics, but to reinterpret it through a minimal SST field theory while preserving its tested limits. The required observable layer is therefore retained as
\begin{equation}
\label{eq:state-labels}
\text{electron state}=(n,\ell,m_\ell,\sigma),
\qquad
\sigma\in\{\!\boldsymbol{\circlearrowleft},\!\boldsymbol{\circlearrowright}\},
\end{equation}
with the intended low-energy identification that the SST branch label $\sigma$ reproduces the observed two-valued internal sector ordinarily encoded by $m_s=\pm \tfrac12$.

A minimally canonical SST atomic equation must satisfy four constraints:
\begin{enumerate}[label=(C\arabic*)]
    \item recover the hydrogenic spectrum in the one-electron weak-coupling limit;
    \item preserve the shell-state counting $N_{n\ell}=2(2\ell+1)$ and $N_n=2n^2$;
    \item admit a closed many-electron mean-field interpretation via a self-consistent density $\rho$;
    \item isolate new SST content to explicit, testable couplings rather than hidden redefinitions.
\end{enumerate}

\section{Axiomatic Foundations and Fluid-Mechanical Ingredients}
The objective is not to replace the successful mathematics of atomic physics, but to derive its structural features---specifically the discrete energy spectrum and the Pauli exclusion geometry---from a minimal, deterministic fluid field theory. To make this manuscript strictly self-contained, we establish the framework upon a continuous, incompressible, inviscid background medium (the condensate), characterized entirely by three primitive constants: an effective mass density $\rho_{\!f}$, a fundamental core radius $r_c$, and a characteristic internal wave speed $\lVert \mathbf{v}_{\!\boldsymbol{\circlearrowleft}}\rVert$.

From this fluid substrate, three functional ingredients are constructed.

\paragraph{(i) Propagating transverse torsion.} Electromagnetic-like radiation is modeled not as point particles, but as transverse director/shear waves in the continuum. Using Cartan's geometric framework for teleparallelism \cite{HehlObukhov2007}, an electromagnetic-strength two-form can be identified as a projection of torsion,
\begin{equation}
\label{eq:Ftorsion}
F \equiv \kappa\,u_a T^a.
\end{equation}
We model the propagating radiation sector as a transverse director field with velocity
\begin{equation}
\label{eq:uA}
\mathbf{u}=\partial_t \mathbf{A},
\qquad
\nabla\cdot \mathbf{A}=0.
\end{equation}
The vector potential $\mathbf{A}$ acts as a spatial displacement field carrying dimensions of length.

\paragraph{(ii) Rotating-fluid source law.} The generation of local circulation from linear forcing is governed by the standard rotating Euler equations \cite{Greenspan1968, Batchelor1967}. For axisymmetric perturbations in a uniformly rotating fluid volume, the axial velocity $w = \partial_t \xi$ sources vertical vorticity $\omega_z$:
\begin{equation}
\label{eq:wz}
\partial_t \omega_z = 2\Omega\,\partial_z w,
\end{equation}
which integrates from rest to
\begin{equation}
\label{eq:wzxi}
\omega_z(r,z,t)=2\Omega\,\partial_z \xi(r,z,t).
\end{equation}
The azimuthal swirl velocity follows kinematically from Stokes' theorem,
\begin{equation}
\label{eq:utheta}
\omega_z = \frac{1}{r}\partial_r(r u_\theta)
\quad\Longrightarrow\quad
u_\theta(r,z,t)=\frac{1}{r}\int_0^r \omega_z(r',z,t)\,r'\,dr'.
\end{equation}
This yields a reversible leading-order swirl response for symmetric up-down forcing, converting macroscopic mechanical displacement into localized vortex circulation.

\paragraph{(iii) Circulation quantization and the Swirl-Clock.}
Following the topological vortex legacy of Kelvin and Helmholtz \cite{Kelvin1869, Helmholtz1858}, we elevate quantized circulation to a primitive invariant of the medium. The fundamental circulation quantum is defined by the core boundary parameters:
\begin{equation}
\label{eq:Gamma0}
\boxed{
\Gamma_0 \equiv 2\pi \rc\,\lVert \mathbf{v}_{\!\boldsymbol{\circlearrowleft}}\rVert,
\qquad
\Gamma = \oint \mathbf v\cdot d\boldsymbol\ell = n\,\Gamma_0, \quad n \in \mathbb{Z}.
}
\end{equation}
To recover relativistic kinematics without a Lorentzian spacetime manifold, we define a scalar bookkeeping field---the Swirl-Clock---which tracks the local time-dilation induced by the fluid's kinetic energy density. We cleanly distinguish the local material turnover scale,
\begin{equation}
\label{eq:Omega0}
\Omega_0 \equiv \frac{\lVert \mathbf{v}_{\!\boldsymbol{\circlearrowleft}}\rVert}{\rc},
\end{equation}
from the laboratory frame, and assign the local clock factors
\begin{equation}
\label{eq:SwirlClock}
\mathbf{S}_{(t)}^{\!\boldsymbol{\circlearrowleft}} = \sqrt{1-\frac{v_{\!\boldsymbol{\circlearrowleft}}^{2}}{c^2}},
\qquad
\mathbf{S}_{(t)}^{\!\boldsymbol{\circlearrowright}} = \sqrt{1-\frac{v_{\!\boldsymbol{\circlearrowright}}^{2}}{c^2}}.
\end{equation}
The labels $(\!\boldsymbol{\circlearrowleft}, \!\boldsymbol{\circlearrowright})$ track the two internal chiral branches (handedness) of the vortex knots, fundamentally reproducing the two-valued internal sector ordinarily encoded by fermionic spin $m_s=\pm \tfrac12$.

\section{Minimal retained field sector}
The first retained sector is a transverse torsion/shear field. The minimal Lagrangian is
\begin{equation}
\label{eq:Ltorsion}
\mathcal{L}_{\rm torsion}
=
\frac12\rhoF\,\lVert \partial_t \mathbf{A}\rVert^2
-
\frac12 K\,\lVert \nabla\times\mathbf{A}\rVert^2,
\qquad
\nabla\cdot\mathbf{A}=0.
\end{equation}
Varying with respect to $\mathbf{A}$ gives
\begin{equation}
\label{eq:Awave-pre}
\rhoF\left(\partial_t^2\mathbf{A}+c^2\nabla\times\nabla\times\mathbf{A}\right)=0,
\qquad
c^2\equiv \frac{K}{\rhoF},
\end{equation}
which reduces, under $\nabla\cdot\mathbf{A}=0$, to
\begin{equation}
\label{eq:Awave}
\boxed{
\left(\frac{1}{c^2}\partial_t^2-\nabla^2\right)\mathbf{A}=0.
}
\end{equation}
This is the propagating field retained from the older torsion/photon sector. No magneto-optical or parity-odd vacuum extensions are retained in the present canonical core.

\section{Minimal retained source sector}
The second retained sector is the rotating-medium source law. The atomic theory proposed below does not require that every torsion excitation originate in a macroscopic rotating cylinder, but the rotating-cylinder calculation provides the cleanest explicit production mechanism for an SST circulation amplitude.

The circulation-foundational choice retained in v0.9 is to normalize the source by local circulation, not by local clock shift. Using Eq.~\eqref{eq:utheta}, define the local loop circulation by
\begin{equation}
\label{eq:GammaLoc}
\boxed{
\Gamma_{\rm loc}(r,z,t)
:=
2\pi r\,u_\theta(r,z,t)
=
2\pi\int_0^r \omega_z(r',z,t)\,r'\,dr'.
}
\end{equation}
The canon-level dimensionless circulation driver is then
\begin{equation}
\label{eq:gammaDef}
\boxed{
\gamma(\mathbf r,t)
\equiv
\frac{\Gamma_{\rm loc}(\mathbf r,t)}{\Gamma_0}.
}
\end{equation}
Equivalently, via Eq.~\eqref{eq:wzxi},
\begin{equation}
\label{eq:gammaXi}
\gamma(r,z,t)
=
\frac{4\pi\Omega}{\Gamma_0}
\int_0^r \partial_z \xi(r',z,t)\,r'\,dr'.
\end{equation}
Thus the new source chain is
\begin{equation}
\label{eq:source-chain}
\xi \longrightarrow \omega_z \longrightarrow u_\theta \longrightarrow \Gamma_{\rm loc} \longrightarrow \gamma,
\end{equation}
with the clock sector treated downstream of circulation rather than as the primitive driver.

For a localized helical packet, a practical canon-compatible ansatz is
\begin{equation}
\label{eq:gammaPacket}
\gamma(r,\phi,z,t)
=
A_\gamma
\exp\!\left(-\frac{r^2}{2w_r^2}\right)
\exp\!\left(-\frac{(z-vt)^2}{2w_z^2}\right)
\cos\!\left(m\phi+kz-\omega t\right).
\end{equation}
This profile is retained because it carries finite amplitude, axial propagation, helical winding, and finite support, while remaining directly interpretable as a circulation-bearing packet.

For continuity with earlier versions, the scalar vorticity-normalized driver
\begin{equation}
\label{eq:ThetaDef}
\Theta(\mathbf r,t) \equiv \frac{\omega_z(\mathbf r,t)}{\Omega_0}
\end{equation}
may still be used as an auxiliary bookkeeping field. In the present v0.9 draft, however, it is no longer the primary canonical source variable.

\section{Atomic sector: branch-resolved minimal Hamiltonian}
The atomic bridge model is built by retaining the ordinary Coulomb core and many-electron screening, and placing the novel SST content entirely in a circulation-sensitive branch potential. The working Hamiltonian is
\begin{equation}
\label{eq:master-H}
\boxed{
\left[
-\frac{\hbar^2}{2\mu}\nabla^2
+V_{\rm core}(r)
+V_{\rm screen}[\rho](\mathbf r)
+V_{\rm circ}^{(\sigma)}[\gamma]
\right]\Psi_{i,\sigma}
=\varepsilon_{i,\sigma}\Psi_{i,\sigma},
}
\end{equation}
with
\begin{equation}
\label{eq:Vcore}
V_{\rm core}(r)=-\frac{Ze^2}{4\pi\varepsilon_0 r}
\end{equation}
for the conservative baseline.

A provisional ``freeze v0.9'' removes the free branch couplings in favor of canonical SST scales. First, define the branch-coupling strength by the small dimensionless material ratio
\begin{equation}
\label{eq:eta0}
\boxed{
\eta_0 \equiv \frac{\lVert\mathbf{v}_{\!\boldsymbol{\circlearrowleft}}\rVert}{c} = \frac{\alpha}{2} \approx 3.648676\times 10^{-3}.
}
\end{equation}
Then take the unperturbed atomic baseline to be branch-symmetric,
\begin{equation}
\label{eq:St0freeze}
S_{t,0}(r)\equiv 1,
\end{equation}
and define the two branch sectors by a circulation-driven clock bookkeeping field
\begin{equation}
\label{eq:StBranches}
\boxed{
S_t^{(\sigma)}(\mathbf r,t)=1+\sigma\,\eta_0\,\gamma(\mathbf r,t),
\qquad
\sigma=\pm 1,
}
\end{equation}
where $\sigma=+1$ corresponds to $\mathbf{S}_{(t)}^{\!\boldsymbol{\circlearrowleft}}$ and $\sigma=-1$ to $\mathbf{S}_{(t)}^{\!\boldsymbol{\circlearrowright}}$. The clock is therefore retained, but only as a downstream representation of circulation-bearing dynamics.

Next, freeze the energy scale to the Rydberg binding scale,
\begin{equation}
\label{eq:ERfreeze}
\boxed{
E_R \equiv \frac{\hbar^2}{2\mu (a_0^{\rm SST})^2}
= \frac{\mu e^4}{2(4\pi\varepsilon_0)^2\hbar^2}.
}
\end{equation}
In v0.4 the clock potential also carried a gradient piece proportional to $(a_0^{\rm SST})^2\abs{\nabla S_t^{(\sigma)}}^2$. In the present v0.8 freeze that term is dropped deliberately rather than silently: once circulation is promoted to the primitive atomic driver, the spatial structure previously delegated to the gradient penalty is instead carried by the explicit circulation field $\gamma(\mathbf r,t)$ and by the orbit-coupled anticommutator term introduced below. The frozen starter model therefore keeps only the algebraic clock shift in $V_{\rm clock}^{(\sigma)}$ and treats spatial/chiral structure as belonging to the circulation sector.

The scalar clock-derived contribution is then
\begin{equation}
\label{eq:Vclock}
\boxed{
V_{\rm clock}^{(\sigma)}
=
E_R\Bigl[1-\bigl(S_t^{(\sigma)}\bigr)^2\Bigr].
}
\end{equation}
At first order this gives
\begin{equation}
\label{eq:Vclock-linear}
\boxed{
\delta V_{\rm clock}^{(\sigma)}
\approx
-2\sigma E_R\eta_0\,\gamma.
}
\end{equation}

The genuinely circulation-foundational addition retained in v0.9 is an orbit-coupled term. To avoid introducing a new fitted constant, the same small canonical factor $\eta_0$ is reused and the coupling is frozen at the same atomic energy scale. The minimal Hermitian choice is the anticommutator form
\begin{equation}
\label{eq:HintCirc}
\boxed{
\delta H_{\rm circ}^{(\sigma)}
=
\frac{E_R\eta_0\,\sigma}{2}
\left\{\frac{\hat L_z}{\hbar},\gamma\right\}.
}
\end{equation}
This form is algebraically equivalent to the previous $\{\hat L_z,\gamma\}/(2\hbar)$ representation, but makes the scaling explicit: $\hat L_z/\hbar$ is dimensionless, so the anticommutator is dimensionless before multiplication by the common frozen atomic energy scale $E_R\eta_0$. The anticommutator is important: for an initial $s$-state with $\hat L_z\psi_{1s}=0$, the operator still yields a nonzero contribution because $\hat L_z$ can act on the helical packet factor inside $\gamma$.

The full frozen linear interaction is therefore
\begin{equation}
\label{eq:HintTotal}
\boxed{
\delta H_{\rm int}^{(\sigma)}
=
-2\sigma E_R\eta_0\,\gamma
+
\frac{E_R\eta_0\,\sigma}{2}
\left\{\frac{\hat L_z}{\hbar},\gamma\right\}.
}
\end{equation}
This is the main circulation-foundational upgrade relative to the earlier scalar-clock baseline: the atomic sector now responds directly to handed circulation, not only to a scalar branch shift.

\subsection{Radiative Stability and the Kelvin-Mode Gap}
A persistent structural objection to classical orbital models is the Larmor formula, which predicts continuous radiation from accelerating charges. In the present fluid framework, continuous dissipation is obstructed by the internal stiffness of the circulation core. The vortex filament supports transverse acoustic oscillations (Kelvin waves) with a fundamental topological excitation gap bounded by the core scale:
\begin{equation}
\label{eq:KelvinGap}
\Delta K \sim \hbar \Omega_0 = \hbar \frac{\lVert \mathbf{v}_{\!\boldsymbol{\circlearrowleft}}\rVert}{\rc} \sim \mathcal{O}(10^2 - 10^3)\ \mathrm{eV}.
\end{equation}
Because the resonant atomic transition energies ($\hbar\omega \approx 1\text{--}10\ \mathrm{eV}$) are orders of magnitude below this internal structural gap, the transverse Kelvin modes of the trefoil core remain frozen. The electron filament therefore acts as a rigid, non-dissipative topological obstruction to the fluid background. The system cannot radiate continuously; energy transfer is strictly gated by the resonant exchange of discrete torsion packets matching the orbital turnover frequencies. This scale separation rigorously justifies truncating the continuous fluid dynamics to the frozen, branch-resolved Hamiltonian of Eq.~\eqref{eq:master-H}.

\section{Many-electron closure and screening}
To reach shell structure, a many-electron mean-field closure is required. The electron density is defined by
\begin{equation}
\label{eq:rho-def}
\rho(\mathbf r)=\sum_{i,\sigma} f_{i,\sigma}\,\abs{\Psi_{i,\sigma}(\mathbf r)}^2,
\end{equation}
with occupation numbers $f_{i,\sigma}\in\{0,1\}$ at the effective level. The minimal retained screening functional is the Hartree term \cite{Hartree1928},
\begin{equation}
\label{eq:Vscreen}
\boxed{
V_{\rm screen}[\rho](\mathbf r)
=
\frac{e^2}{4\pi\varepsilon_0}
\int \frac{\rho(\mathbf r')}{\abs{\mathbf r-\mathbf r'}}\,d^3r'.
}
\end{equation}
No SST-specific correction to screening is introduced in the canonical starter model. This is deliberate: a predictive theory must first recover the ordinary many-electron limit before attempting branch- or medium-induced corrections to screening.

At the level of the working atomic Hamiltonian, the shell-state capacity follows the standard counts:
\begin{equation}
\label{eq:subshell-count}
N_{n\ell}=2(2\ell+1),
\end{equation}
and
\begin{equation}
\label{eq:shell-count}
N_n=\sum_{\ell=0}^{n-1}2(2\ell+1)=2n^2.
\end{equation}
In the present bridge model, the factor $2\ell+1$ arises from the ordinary rotational sectors $m_\ell=-\ell,\ldots,+\ell$, while the factor $2$ is interpreted as the low-energy branch doublet associated with $\mathbf{S}_{(t)}^{\!\boldsymbol{\circlearrowleft}}$ and $\mathbf{S}_{(t)}^{\!\boldsymbol{\circlearrowright}}$.

Crucially, this exclusion is not a fundamental postulate of SST, but an emergent hydrodynamic necessity. As proven in Appendix~\ref{app:PauliBarrier} and formalized in Section~\ref{sec:MultiElectronShells}, forcing two identically handed trefoils into the same spatial domain generates a localized dynamic pressure spike that exceeds the fluid's cavitation limit. This topological repulsion acts as an absolute physical barrier, rigorously deriving the Pauli exclusion principle and the Aufbau architecture strictly from fluid geometry.

\section{Hydrogenic limit and SST length scale}
The model must recover the ordinary hydrogenic spectrum when
\begin{equation}
V_{\rm screen}\to 0,
\qquad
\gamma\to 0,
\qquad
S_t^{(+)}=S_t^{(-)}.
\end{equation}
Then~\eqref{eq:master-H} reduces to the standard one-electron Coulomb problem,
\begin{equation}
\label{eq:En}
E_n=-\frac{\mu Z^2 e^4}{2(4\pi\varepsilon_0)^2\hbar^2}\frac{1}{n^2}.
\end{equation}
The SST atomic scale is tied to the canonical constants by
\begin{equation}
\label{eq:a0-sst}
\boxed{
a_0^{\rm SST}
=
\frac{F_{\rm swirl}^{\max}\,\rc^2}{M_e\,\lVert\mathbf{v}_{\!\boldsymbol{\circlearrowleft}}\rVert^2}
=
\frac{2\rc}{\alpha^2}.
}
\end{equation}
Using the canonical values
\begin{equation}
\lVert\mathbf{v}_{\!\boldsymbol{\circlearrowleft}}\rVert = 1.09384563\times 10^6\ \mathrm{m\,s^{-1}},
\quad
\rc = 1.40897017\times 10^{-15}\ \mathrm{m},
\quad
F_{\rm swirl}^{\max}=29.053507\ \mathrm{N},
\end{equation}
this gives
\begin{equation}
\label{eq:a0-num}
a_0^{\rm SST}
=
5.29177161\times 10^{-11}\ \mathrm{m},
\end{equation}
while the equivalent relation $2\rc/\alpha^2$ yields
\begin{equation}
\label{eq:a0-num2}
\frac{2\rc}{\alpha^2}=5.29177214\times 10^{-11}\ \mathrm{m},
\end{equation}
consistent with the ordinary Bohr scale to the expected rounding level.

This gives the hydrogenic SST radius law
\begin{equation}
\label{eq:RnZ}
R_{n,Z}^{\rm SST}\sim \frac{n^2}{Z}\,a_0^{\rm SST}.
\end{equation}
Thus the atomic bridge model retains the tested size law while reinterpreting $a_0$ as a medium-derived coherence scale. In the present draft this agreement is to be read as a \emph{consistency check on the retained canonical scales}, not as an independent prediction, since $(F_{\rm swirl}^{\max},\rc,\lVert\mathbf{v}_{\!\boldsymbol{\circlearrowleft}}\rVert)$ were fixed upstream in the SST program.

\paragraph{Thermodynamic Equilibrium of the Ground State.}
The scaling in Eq.~\eqref{eq:a0-sst} is not merely a dimensional coincidence; it reflects a stationary hydrodynamic equilibrium. In the surrounding swirl condensate, the unperturbed ground state corresponds to a zero effective ``swirl temperature'' ($T_{\rm swirl} = 0$), where geometric strain vanishes.

The stable orbital radius $a_0^{\rm SST}$ emerges precisely at the surface where the outward dynamic swirl pressure balances the inward coherent vacuum tension. Treating the core as a source of persistent circulation, the local dynamic pressure scales as $P_{\rm dyn} \sim \frac{1}{2}\rhoF v_\theta^2$. Enforcing a steady-state Euler balance between the core circulation and the far-field envelope yields the ratio of the core radius to the equilibrium boundary:
\begin{equation}
\label{eq:PressureBalance}
\frac{a_0^{\rm SST}}{\rc} = \frac{2c^2}{\lVert \mathbf{v}_{\!\boldsymbol{\circlearrowleft}}\rVert^2} = \frac{2}{\eta_0^2} = \frac{8}{\alpha^2}.
\end{equation}
Thus, the hydrogenic orbital is reinterpreted not as a probabilistic cloud, but as a sharply defined, thermodynamically stable flow boundary within the incompressible medium. Transitioning between orbitals requires discrete injections of SST Heat (Kelvin-mode acoustic energy) to deform this equilibrium boundary against the fluid pressure.

\section{Hydrodynamic Derivation of Multi-Electron Shells and the Aufbau Principle}
\label{sec:MultiElectronShells}

In the standard phenomenological bridge (Sec.~6), the many-electron atomic structure is governed by an effective Hartree screening potential and the posited Pauli state-counting rule $N_n = 2n^2$. To render the Swirl-String Theory canonically closed, these features must be derived directly from the continuum mechanics of the swirl condensate and the delay-induced mode selection of the NLS framework.

\subsection{Screened Swirl Circulation and the Effective Euler Balance}
Consider a central nucleus of topological charge $Z$. In the ideal fluid limit, this acts as a macroscopic circulation sink/source, establishing a baseline azimuthal swirl field $\mathbf{v}_{\theta}^{\rm nuc}(r) = Z \Gamma_0 / (2\pi r) \hat{\boldsymbol{\theta}}$.

As $N$ electron trefoils are added to the atomic basin, they act as localized, screened circulation defects. The macroscopic steady-state fluid velocity at a radial distance $r$ is determined by the total enclosed circulation via Stokes' theorem:
\begin{equation}
\label{eq:GammaEff}
\Gamma_{\rm eff}(r) = \oint_{\partial S_r} \mathbf{v}_{\theta} \cdot d\mathbf{l} = \Gamma_0 \left( Z - \int_0^r \rho_\sigma(\mathbf{r}') \, d^3r' \right) \equiv \Gamma_0 Z_{\rm eff}(r),
\end{equation}
where $\rho_\sigma(\mathbf{r}')$ is the number density of electron trefoils. The resulting outward dynamic pressure of the macroscopic atomic vortex is governed by the Euler equation for steady, incompressible flow:
\begin{equation}
\label{eq:EulerDynamicPressure}
\nabla P_{\rm dyn} = \rho_{\!f} \left( \mathbf{v}_{\theta} \cdot \nabla \right) \mathbf{v}_{\theta} \quad \Longrightarrow \quad P_{\rm dyn}(r) = \frac{1}{2}\rho_{\!f} \left( \frac{\Gamma_{\rm eff}(r)}{2\pi r} \right)^2.
\end{equation}

\subsection{Delay-Induced Orbital Quantization}
In the ideal fluid limit, a coherent localized excitation is governed by the Gross-Pitaevskii or Nonlinear Schr\"odinger (NLS) equation for the macroscopic wavefunction of the condensate \cite{Pitaevskii2003}. A macroscopic electron shell corresponds to a coherent, phase-locked soliton envelope within the atomic basin.

Closed circulation loops with finite propagation speeds constitute a generic class of delayed feedback systems. As established in the dynamics of delay-differential equations (DDEs), temporal feedback delays alone are sufficient to enforce spectral discreteness through dynamical stability \cite{Erneux2009}. The orbital transit time for a fluid element at radius $r_n$ is:
\begin{equation}
\label{eq:OrbitalDelay}
\tau_{\rm orb} = \frac{2\pi r_n}{v_{\theta}(r_n)} = \frac{(4\pi^2 r_n^2)}{\Gamma_0 Z_{\rm eff}(r_n)}.
\end{equation}
For the orbital to be thermodynamically stable against radiative decay (Kelvin-mode dissipation), this macroscopic delay must rationally lock with the intrinsic temporal periodicity of the electron trefoil core, $\tau_c = 2\pi r_c / \lVert\mathbf{v}_{\!\boldsymbol{\circlearrowleft}}\rVert$. This is the fluid-mechanical equivalent of the Adler phase-locking equation with delay \cite{Adler1946}. The resonant locking condition forces the continuous spectrum to collapse into discrete bands; for a mode of order $n$ the orbital action scales quadratically with the integer winding index $n \in \mathbb{Z}^+$:
\begin{equation}
\label{eq:DelayLocking}
r_n v_{\theta}(r_n) = n \left( r_c \lVert\mathbf{v}_{\!\boldsymbol{\circlearrowleft}}\rVert \right) \left( \frac{2}{\eta_0^2} \right) \frac{1}{Z_{\rm eff}(r_n)}.
\end{equation}
Substituting the canonical definition of the SST Bohr radius $a_0^{\rm SST} = 2r_c / \eta_0^2$ (Eq.~\eqref{eq:a0-sst}) directly recovers the exact radial quantization of the atomic shells:
\begin{equation}
\label{eq:SST_OrbitRadii}
\boxed{
r_n = \frac{n^2}{Z_{\rm eff}(r_n)} \, a_0^{\rm SST}.
}
\end{equation}
The hydrogenic spectrum is thereby derived not from a probability wave, but as the discrete set of radii where the macroscopic circulation delay perfectly phase-locks with the topological core of the constituent $2{:}3$ trefoils.

\subsection{Topological Derivation of the $2n^2$ Shell Capacity}
The most profound consequence of replacing abstract quantum operators with hydrodynamic topology is the derivation of the Aufbau shell limits. Why can shell $n$ hold exactly $N_n = 2n^2$ electrons?

A macroscopic shell of index $n$ supports orthogonal fluid-acoustic standing waves (spherical harmonics $Y_{\ell m_\ell}$) whose total mode count up to degree $n-1$ is strictly fixed by 3D spherical geometry:
\begin{equation}
\label{eq:SphericalHarmonics}
\sum_{\ell=0}^{n-1} (2\ell + 1) = n^2.
\end{equation}
Each of these $n^2$ spatial acoustic nodes provides a localized minimum in the dynamic pressure envelope where an electron trefoil can stably reside.

We must now invoke the absolute Cavitation Limit (the topological Pauli barrier derived in Appendix~\ref{app:PauliBarrier}). If two identical trefoils are forced into the same spatial pressure node, their Biot-Savart interaction energy induces a local pressure drop $\Delta P \sim -\frac{1}{2}\rho_{\!f} \lVert 2\mathbf{v}_{\!\boldsymbol{\circlearrowleft}} \rVert^2$ that exceeds the universal cavitation threshold $F_{\rm swirl}^{\max} \approx 29.05\ \mathrm{N}$. The condensate physically ruptures, forbidding the state.

However, if the two trefoils possess opposite Swirl-Clock internal branches ($\sigma_1 = +1$ and $\sigma_2 = -1$), their circulation vectors at the boundary interface are strictly anti-parallel ($d\mathbf{l}_1 \cdot d\mathbf{l}_2 < 0$). This destructive interference locally \emph{lowers} the kinetic energy density, avoiding the cavitation limit and creating a stable, bound pair (a local geometric Cooper-like pairing of opposite chiralities).

Therefore, every spatial acoustic node can support exactly two trefoils of opposite internal handedness before the fluid physically cavitates. The maximum electron capacity of the $n$-th hydrodynamic shell is exactly:
\begin{equation}
\label{eq:AufbauDerived}
\boxed{
N_n = 2 \times (\text{spatial fluid nodes}) = 2n^2.
}
\end{equation}
The Pauli exclusion principle and the architecture of the periodic table are thus derived entirely from the finite strength of the vacuum fluid ($F_{\rm swirl}^{\max}$) and the discrete harmonic geometry of 3D incompressible flow.

\subsection{High-$Z$ Core Contraction and Swirl-Clock Kinematics}
To encompass all known elements across the periodic table, the theory must also account for the relativistic kinematics of inner-shell electrons in heavy atoms. For high-$Z$ nuclei, the macroscopic fluid velocity near the nucleus $v_{\theta}^{\rm nuc} = Z \Gamma_0 / (2\pi r)$ becomes a significant fraction of the canonical medium speed $\lVert\mathbf{v}_{\!\boldsymbol{\circlearrowleft}}\rVert$. In this regime, the static Swirl-Clock field $S_t^{(\sigma)}(\mathbf r)$ strongly deviates from unity.

The kinetic operator in the effective Hamiltonian (Eq.~\eqref{eq:master-H}) must be covariant with respect to the local clock foliation, modifying the effective inertial mass:
\begin{equation}
\label{eq:Hkin_HighZ}
-\frac{\hbar^2}{2\mu}\nabla^2 \quad \longrightarrow \quad -\frac{\hbar^2}{2\mu} \nabla \cdot \left( \bigl(S_t^{(\sigma)}(\mathbf r)\bigr)^2 \nabla \right).
\end{equation}
Because $S_t < 1$ deep in the nuclear swirl basin, the effective radial kinetic energy is suppressed, requiring the inner orbitals to contract further inward to maintain the Euler pressure balance of Eq.~\eqref{eq:EulerDynamicPressure}. This directly recovers the relativistic $1s$ orbital shrinkage responsible for the macroscopic chemical properties of heavy metals (e.g., the inertness of Gold), derived here entirely from classical fluid time-dilation rather than Dirac spinor kinematics.

\section{Transition amplitudes and torsion capture}
A circulation-bearing torsion packet~\eqref{eq:gammaPacket} drives atomic transitions through the frozen interaction~\eqref{eq:HintTotal}. For an initial state $\ket{i,\sigma}$ and final state $\ket{f,\sigma'}$, the first-order matrix element is
\begin{equation}
\label{eq:Mfi}
M_{fi}^{\sigma'\sigma}=\matrixel{f,\sigma'}{\delta H_{\rm int}}{i,\sigma}.
\end{equation}
The conservative resonance condition is
\begin{equation}
\label{eq:resonance}
\hbar\omega \approx \varepsilon_f-\varepsilon_i.
\end{equation}
If the packet carries azimuthal phase factor $e^{im\phi}$, then the orbital projection rule remains
\begin{equation}
\label{eq:dm}
\Delta m_\ell = m.
\end{equation}
The new feature is that the anticommutator term in Eq.~\eqref{eq:HintCirc} is explicitly odd under $m\mapsto -m$. For the helical packet ansatz of Eq.~\eqref{eq:gammaPacket}, the radial and axial Gaussian envelopes are $\phi$-independent, so $\hat L_z=-i\hbar\partial_\phi$ acts only on the azimuthal factor. One therefore has the exact identity
\begin{equation}
\label{eq:Lzgamma}
\hat L_z\,\gamma = m\hbar\,\gamma,
\end{equation}
so the interaction can distinguish opposite-handed packets at operator level rather than only after selecting final states. This is exactly the chiral handle absent from the scalar-only v0.4 freeze.

\paragraph{Operator-level $1s$ asymmetry remark.}
For the hydrogenic ground state, $\hat L_z\psi_{1s}=0$. Applying Eq.~\eqref{eq:HintTotal} together with Eq.~\eqref{eq:Lzgamma} gives
\begin{equation}
\label{eq:Hint1s}
\boxed{
\delta H_{\rm int}^{(\sigma)}\psi_{1s}
=
E_R\eta_0\,\sigma\left(-2+\frac{m}{2}\right)\gamma\,\psi_{1s}.
}
\end{equation}
Thus the helical packet channels are already inequivalent before continuum integration:
\begin{equation}
\label{eq:Hint1sPm}
 m=+1:\quad \delta H_{\rm int}^{(\sigma)}\psi_{1s}=-\frac{3}{2}E_R\eta_0\,\sigma\,\gamma\,\psi_{1s},
 \qquad
 m=-1:\quad \delta H_{\rm int}^{(\sigma)}\psi_{1s}=-\frac{5}{2}E_R\eta_0\,\sigma\,\gamma\,\psi_{1s}.
\end{equation}
If the corresponding continuum overlap integrals satisfy $|I_{+1}(k,\omega)|=|I_{-1}(k,\omega)|$ to leading order, then the source-level rate asymmetry is already nonzero. For the Coulomb baseline this equality is the natural leading-order expectation: the bound--continuum matrix element factorizes into radial and angular pieces, the radial integral depends on $k$ and the radial profile of $\gamma$ but not on the sign of $m$, and the sign of $m$ enters only through the azimuthal selection rule $\Delta m_\ell=m$.
\begin{equation}
\label{eq:A1sSource}
\boxed{
\mathcal A_{1s}^{\rm(source)}
=
\frac{|{-3}/{2}|^2-|{-5}/{2}|^2}{|{-3}/{2}|^2+|{-5}/{2}|^2}
=
-\frac{8}{17}
\approx -0.4706.
}
\end{equation}
Equation~\eqref{eq:A1sSource} is not yet the full continuum photoionization result; it is the operator-level asymmetry implied by the frozen interaction before the Coulomb-continuum matrix elements are evaluated. The negative sign means that, at this source level, the $h=-1$ packet channel is favored over the $h=+1$ channel because the prefactor magnitude $|{-5}/{2}|$ exceeds $|{-3}/{2}|$.

This suggests the following interpretation. Upward shell excitation under resonant incoming torsion is the atomic analogue of a ``giant jet,'' while downward relaxation accompanied by outgoing torsion radiation is the analogue of a ``sprite.'' These are analogical labels only; the retained mathematics is ordinary time-dependent perturbation theory acting on a branch-resolved Hamiltonian.

\section{Continuum sector and photoionization benchmark}
The sharpest next benchmark is photoionization rather than further shell bookkeeping. A branch-resolved atomic theory is incomplete until it specifies not only its bound spectrum but also its continuum and the dynamical channel by which an incoming torsion packet transfers energy into that continuum. The minimal conservative choice is to define the continuum as the positive-energy sector of the same branch-resolved Hamiltonian already introduced in Eq.~\eqref{eq:master-H}.

Because $S_{t,0}\equiv 1$ in the present freeze, the static clock contribution vanishes identically. The unperturbed branch Hamiltonian is therefore simply
\begin{equation}
\label{eq:H0sigma}
\boxed{
H_0^{(\sigma)}
=
-\frac{\hbar^2}{2\mu}\nabla^2
+
V_{\rm core}(r)
+
V_{\rm screen}[\rho](\mathbf r),
\qquad
\sigma=\pm 1.
}
\end{equation}
The corresponding effective potential is
\begin{equation}
\label{eq:Veffsigma}
V_{\rm eff}^{(\sigma)}(\mathbf r)
=
V_{\rm core}(r)
+
V_{\rm screen}[\rho](\mathbf r).
\end{equation}
Bound states satisfy
\begin{equation}
\label{eq:boundstates}
H_0^{(\sigma)}\psi_{n\ell m_\ell}^{(\sigma)}
=
E_{n\ell}^{(\sigma)}\psi_{n\ell m_\ell}^{(\sigma)},
\qquad
E_{n\ell}^{(\sigma)}<0,
\end{equation}
whereas continuum scattering states satisfy
\begin{equation}
\label{eq:continuumstates}
H_0^{(\sigma)}\psi_{\mathbf k}^{(\sigma)}
=
E_{\mathbf k}^{(\sigma)}\psi_{\mathbf k}^{(\sigma)},
\qquad
E_{\mathbf k}^{(\sigma)}>0.
\end{equation}
Under the localized-packet assumption
\begin{equation}
\label{eq:gammaAsymptotic}
\gamma(\mathbf r,t)\to 0 \qquad (\lVert \mathbf r\rVert\to\infty),
\end{equation}
we have $S_t^{(\sigma)}\to 1$ and hence no asymptotic branch-dependent offset. The continuum threshold is therefore branch-degenerate and may be taken as
\begin{equation}
\label{eq:continuum-threshold}
E_{\rm cont}^{(+)}(0)=E_{\rm cont}^{(-)}(0)=0.
\end{equation}

The next required step is to close the dynamical link between the propagating torsion field and the circulation driver $\gamma$. The retained v0.8 closure is to extract vorticity from the field velocity $\mathbf u=\partial_t\mathbf A$ and then build the loop circulation from that vorticity. Define
\begin{equation}
\label{eq:omegaA}
\omega_z^{(A)}(\mathbf r,t)
:=
\hat{\mathbf z}\cdot\left(\nabla\times \partial_t \mathbf A\right)(\mathbf r,t),
\end{equation}
then
\begin{equation}
\label{eq:GammaA}
\Gamma_{\rm loc}^{(A)}(r,z,t)
:=
2\pi\int_0^r \omega_z^{(A)}(r',z,t)\,r'\,dr',
\end{equation}
and finally
\begin{equation}
\label{eq:gammaA}
\boxed{
\gamma^{(A)}(\mathbf r,t)
=
\frac{\Gamma_{\rm loc}^{(A)}(\mathbf r,t)}{\Gamma_0}.
}
\end{equation}
This is the key circulation-foundational replacement for the older scalar-insertion strategy. The field now drives the atom through canon-normalized circulation.

To permit back-reaction one must source the field equation itself. The branch-sensitive sourced wave equation is retained as
\begin{equation}
\label{eq:Awave-source}
\boxed{
\left(\frac{1}{c^2}\partial_t^2-\nabla^2\right)\mathbf A
=
g_J \,\widetilde{\mathbf J}_{\rm br},
}
\end{equation}
with the source length frozen to the canonical core scale
\begin{equation}
\label{eq:gJfreeze}
\boxed{
 g_J \equiv r_c = \frac{\alpha^2}{2}\,a_0^{\rm SST}.
}
\end{equation}
The corresponding branch current is
\begin{equation}
\label{eq:Jbr}
\boxed{
\mathbf J_{\rm br}
=
\frac{\hbar}{2 i \mu}
\sum_{\sigma=\pm 1}
\sigma
\left(
\Psi_\sigma^*\nabla\Psi_\sigma
-
(\nabla\Psi_\sigma^*)\Psi_\sigma
\right),
\qquad
\widetilde{\mathbf J}_{\rm br}\equiv \frac{\mathbf J_{\rm br}}{\Omega_0}.
}
\end{equation}
Equation~\eqref{eq:Jbr} remains a provisional closure ansatz rather than a theorem of SST, but it keeps the field-atom-source loop closed with no free prefactor.

Once $\gamma$ is defined by Eq.~\eqref{eq:gammaA}, the frozen circulation interaction from Eq.~\eqref{eq:HintTotal} becomes the genuine photoionization operator. The branch-resolved ionization threshold from an initial bound state $i$ to a final continuum branch $\sigma_f$ is
\begin{equation}
\label{eq:IonizationThresholdGeneral}
I_i^{(\sigma_i\to \sigma_f)}
=
E_{\rm cont}^{(\sigma_f)}(0)-E_i^{(\sigma_i)}.
\end{equation}
Under the asymptotic choice~\eqref{eq:continuum-threshold}, this reduces to
\begin{equation}
\label{eq:IonizationThresholdSimple}
I_i^{(\sigma_i)}=-E_i^{(\sigma_i)}.
\end{equation}
The photoionization condition is then frequency-thresholded,
\begin{equation}
\label{eq:PhotoThreshold}
\boxed{
\hbar\omega \ge I_i^{(\sigma_i\to \sigma_f)}.
}
\end{equation}
For an emitted free electron with asymptotic kinetic energy $K$, one obtains
\begin{equation}
\label{eq:KineticEnergyAtom}
\boxed{
K
=
\hbar\omega - I_i^{(\sigma_i\to \sigma_f)}.
}
\end{equation}
For a condensed system with work function $\phi^{(\sigma_i\to \sigma_f)}$, the analogous relation is
\begin{equation}
\label{eq:PhotoelectricSolid}
\boxed{
K_{\max}
=
\hbar\omega - \phi^{(\sigma_i\to \sigma_f)}.
}
\end{equation}
Thus the threshold is controlled by frequency, while the intensity enters only through the occupation number of the incoming torsion mode.

To show this explicitly, let $N_\omega$ denote the occupation number of an incoming torsion packet of frequency $\omega$. Then the continuous version of Fermi's golden rule \cite{Dirac1927} becomes
\begin{equation}
\label{eq:GoldenRuleContinuum}
\boxed{
\Gamma_{i\to \mathbf k,\sigma_f}
=
\frac{2\pi}{\hbar}
N_\omega
\left|
\matrixel{\psi_{\mathbf k}^{(\sigma_f)}}{\delta H_{\rm int}^{(\sigma_f)}}{\psi_i^{(\sigma_i)}}
\right|^2
\delta\!\left(
E_{\mathbf k}^{(\sigma_f)}-E_i^{(\sigma_i)}-\hbar\omega
\right).
}
\end{equation}
This is the required rate structure if the SST torsion packet is to reproduce the key empirical logic of the photoelectric effect: an immediate threshold in frequency, and an emission rate proportional to flux or intensity rather than a threshold controlled by amplitude.

\textbf{Analogy suitable for a smart teenager:} Think of the atom as a perfectly tuned bell, and the incoming torsion wave as a precise sequence of taps. If the taps arrive at the right speed (frequency) and match the bell's internal twisting direction (handedness), the bell rings and ejects an electron. If the twisting is backwards, the taps just glance off and fail to excite it.

A useful distinction therefore emerges between two regimes. In the \emph{weak-field quantum absorption regime}, Eqs.~\eqref{eq:PhotoThreshold}--\eqref{eq:GoldenRuleContinuum} define the canonical baseline. In a possible later \emph{strong-field nonlinear regime}, sufficiently large coherent $\gamma$ could reshape the effective barrier and produce field-emission or over-the-barrier escape. That second regime should be treated as an extension and not conflated with the baseline photoelectric benchmark.

The branch-sensitive content now enters through an explicitly handed operator. Define the handedness observable
\begin{equation}
\label{eq:Asymmetry}
\boxed{
\mathcal A(\omega)
=
\frac{\Gamma_{h=+1}-\Gamma_{h=-1}}{\Gamma_{h=+1}+\Gamma_{h=-1}},
}
\end{equation}
where $h=\pm 1$ labels the handedness of the incoming torsion packet. For the helical packet ansatz~\eqref{eq:gammaPacket} at fixed propagation direction, the handedness is identified with the sign of the azimuthal winding number,
\begin{equation}
\label{eq:h-sign-m}
\boxed{h\equiv \mathrm{sgn}(m).}
\end{equation}
Combined with the orbital rule $\Delta m_\ell=m$ from Eq.~\eqref{eq:dm}, the anticommutator term in Eq.~\eqref{eq:HintCirc} implies that same-handed channels acquire an operator-level factor proportional to $m$. The corresponding starter-model overlap preference is
\begin{equation}
\label{eq:handed-selection}
 h=+1:\ \sigma_i=\sigma_f=+1\  \text{favored},
 \qquad
 h=-1:\ \sigma_i=\sigma_f=-1\  \text{favored},
\end{equation}
with opposite-handed channels suppressed by reduced helical overlap and branch-flip channels absent at linear order. In contrast to v0.4, this preference is now traceable directly to the interaction operator rather than only to packet labeling.

If $\mathcal A(\omega)=0$ identically even after including Eq.~\eqref{eq:HintCirc}, then the circulation-foundational extension still collapses to a relabeling. If instead $\mathcal A(\omega)\neq 0$ with a definite symmetry and scale, then the theory acquires a genuinely SST-specific atomic signal. This is the main computational fork opened by the present v0.9 draft.


\section{From a closed circulation ring and a torsion wave to trefoil closure}
\label{sec:ring-to-trefoil}

We now sharpen the circulation-foundational picture behind the frozen interaction by asking what additional condition is required for a closed circulation carrier plus a helical torsion mode to realize a trefoil closure. The primitive SST ingredients remain the circulation and turnover scales
\begin{equation}
\Gamma_0 = 2\pi r_c\,\lVert \mathbf{v}_{\!\boldsymbol{\circlearrowleft}}\rVert,
\qquad
\Omega_0 = \frac{\lVert \mathbf{v}_{\!\boldsymbol{\circlearrowleft}}\rVert}{r_c},
\qquad
f_0 = \frac{\Omega_0}{2\pi},
\label{eq:Gamma0Omega0Trefoil}
\end{equation}
with numerical values
\begin{equation}
\Gamma_0 \approx 9.68361920\times 10^{-9}\ \mathrm{m^2\,s^{-1}},
\qquad
\Omega_0 \approx 7.76344066\times 10^{20}\ \mathrm{s^{-1}},
\qquad
f_0 \approx 1.23558996\times 10^{20}\ \mathrm{Hz}.
\label{eq:Gamma0Omega0TrefoilNum}
\end{equation}
In the present canonical parameter set, \(\Omega_0\) coincides numerically with the electron Compton angular scale.

\paragraph{Step 1: the unknotted closed ring.}
Take a closed circulation ring of major radius \(R\), with centerline parameter \(\theta\in[0,2\pi)\):
\begin{equation}
\mathbf{X}_{\mathrm{ring}}(\theta)
=
\begin{pmatrix}
R\cos\theta\\
R\sin\theta\\
0
\end{pmatrix}.
\label{eq:ringcenterlineTrefoil}
\end{equation}
This is topologically the unknot, and the circulation-preserving ideal-flow reading traces back to Helmholtz and Kelvin \cite{Helmholtz1858,Kelvin1869}.

\paragraph{Step 2: a helical torsion mode on the ring.}
Let a torsion/Kelvin-like mode propagate on the ring with phase
\begin{equation}
\Phi(\theta,t)=q\theta-\omega_\tau t+\phi_0,
\label{eq:poloidalphaseTrefoil}
\end{equation}
where \(q\in\mathbb Z\) is the integer winding number of the torsion modulation and \(\omega_\tau\) is the torsion-wave frequency. A minimal toroidal--poloidal embedding is
\begin{equation}
\mathbf{X}(\theta,t)
=
\begin{pmatrix}
\bigl(R+a\cos\Phi(\theta,t)\bigr)\cos\!\bigl(p\theta-\Omega t\bigr)\\[4pt]
\bigl(R+a\cos\Phi(\theta,t)\bigr)\sin\!\bigl(p\theta-\Omega t\bigr)\\[4pt]
a\sin\Phi(\theta,t)
\end{pmatrix},
\qquad 0<a\ll R,
\label{eq:torusknotansatzTrefoil}
\end{equation}
where \(p\in\mathbb Z\) counts the toroidal winding and \(q\in\mathbb Z\) the poloidal winding.

\paragraph{Step 3: closure condition.}
A torus-knot closure occurs only if the toroidal and poloidal windings close after a common period; this is the standard torus-knot closure condition used across vortex and wave-knot constructions \cite{Kedia2013,BerryDennis2001}:
\begin{equation}
p,q\in\mathbb Z,
\qquad
\gcd(p,q)=1.
\label{eq:torusknotclosureTrefoil}
\end{equation}
The simplest nontrivial knot is the trefoil:
\begin{equation}
T_{2,3}\quad\text{or equivalently}\quad T_{3,2}.
\label{eq:trefoiltypesTrefoil}
\end{equation}
Here \(T_{2,3}\) and \(T_{3,2}\) denote the same knot type; they differ only by which winding is called toroidal and which poloidal.

\paragraph{Step 4: Delay-induced mode selection and phase locking.}
In a closed swirl string, the finite transport speed $\lVert \mathbf{v}_{\!\boldsymbol{\circlearrowleft}}\rVert$ along the geometric circumference $L \approx 2\pi R$ defines a fundamental circulation delay,
\begin{equation}
\label{eq:CirculationDelay}
\tau_c = \frac{L}{\lVert \mathbf{v}_{\!\boldsymbol{\circlearrowleft}}\rVert}.
\end{equation}
When the transverse torsion field propagates on this closed topology, the string acts as a delayed feedback system. The internal state of the excitation is characterized by the phase variable $\Phi(\theta,t)$ from Eq.~\eqref{eq:poloidalphaseTrefoil}, which functions as an intrinsic clock associated with the circulating structure.

Because the torsion wave continuously self-interacts after one round-trip delay $\tau_c$, the temporal delay alone acts as a dynamical selection filter. The transcendental locking condition for a circulating phase oscillator with delay forces the spectral continuous band to collapse into a discrete ladder of stability-selected modes. Trefoil closure therefore requires rational phase locking between the toroidal and poloidal turnover frequencies:
\begin{equation}
\frac{\Omega_{\mathrm{pol}}}{\Omega_{\mathrm{tor}}}
=
\frac{q}{p}
=
\frac{3}{2}
\quad\text{or}\quad
\frac{2}{3},
\label{eq:trefoilphaselockTrefoil}
\end{equation}
depending on which motion is designated toroidal and which poloidal. In the present SST reading, the primitive turnover scale is set by \(\Omega_0\), so the natural delay-induced stability resonance is
\begin{equation}
\Omega_{\mathrm{tor}}\sim \Omega_0,
\qquad
\Omega_{\mathrm{pol}}\sim \frac{3}{2}\Omega_0
\quad\text{or}\quad
\Omega_{\mathrm{tor}}\sim \frac{2}{3}\Omega_{\mathrm{pol}}.
\label{eq:trefoilOmega0Trefoil}
\end{equation}

\begin{equation}
\boxedstatement{The primitive internal rate $\Omega_0 \approx \omega_c$ supplies the natural turnover scale, but discrete trefoil formation is dynamically enforced by the temporal circulation delay $\tau_c$, which stabilizes only the rational $2{:}3$ mode lock.}
\label{eq:boxedOmega0Lock}
\end{equation}

\paragraph{Step 5: circulation-first ordering.}
The local circulation carried by the ring is
\begin{equation}
\Gamma_C(t)=\oint_C \mathbf{u}\cdot d\boldsymbol{\ell},
\label{eq:loopcirculationTrefoil}
\end{equation}
and in the SST ideal limit the primitive circulation unit is \(\Gamma_0\). A minimal quantized closure condition is therefore
\begin{equation}
\Gamma_C = N_\Gamma\,\Gamma_0,
\qquad
N_\Gamma\in\mathbb Z.
\label{eq:circulationquantizationTrefoil}
\end{equation}
The ontological order is then, in a circulation-first spirit consistent with helicity-based knot dynamics \cite{Moffatt1969},
\begin{equation}
\Gamma,\ \boldsymbol{\omega},\ u_\theta
\;\longrightarrow\;
\text{closure/helicity}
\;\longrightarrow\;
S_t^{\boldsymbol{\circlearrowleft}},\ S_t^{\boldsymbol{\circlearrowright}}.
\label{eq:circulationbeforeclockTrefoil}
\end{equation}

\begin{equation}
\boxedstatement{Circulation, vorticity, and swirl are primary; the clock fields are downstream bookkeeping of circulation-bearing structure.}
\label{eq:boxedCirculationBeforeClock}
\end{equation}

\paragraph{Step 6: Why temporal delay is required for quantization.}
A small-amplitude helical disturbance of an unknot generically gives a deformed, radiatively unstable ring, not a nontrivial knot. Spatial boundary conditions are traditionally invoked to enforce mode discretization. However, the closed-loop delay framework demonstrates that temporal circulation delays alone are sufficient to enforce spectral discreteness through dynamical stability. Trefoil formation definitively requires:
\begin{enumerate}
    \item a \emph{closed} circulation carrier defining a strict feedback delay $\tau_c$,
    \item a \emph{rational} toroidal--poloidal phase lock induced by the self-interaction of the torsion field over this delay,
    \item topological persistence against continuous relaxation back to the unknot.
\end{enumerate}
Thus the trefoil must be viewed as a delay-induced, mode-locked closure state, not as a generic kinematic consequence of any torsion wave.

\begin{equation}
\boxedstatement{A torsion wave does not automatically create a trefoil; it does so only if the combined toroidal and poloidal motions lock into the coprime $2{:}3$ or $3{:}2$ closure ratio.}
\label{eq:boxedTrefoilNotAutomatic}
\end{equation}

\paragraph{Step 7: energetic admissibility without reopening free parameters.}
To avoid reopening unfrozen coefficients, we write the energetic admissibility condition in normalized form,
\begin{equation}
E_{\mathrm{eff}}[K]
=
E_R\,\mathcal E_{\mathrm{top}}[K],
\label{eq:EeffTopNormalized}
\end{equation}
where \(E_R\) is the frozen atomic scale already used in the main text and \(\mathcal E_{\mathrm{top}}[K]\) is a dimensionless topological-functional placeholder. In the present working draft, \(\mathcal E_{\mathrm{top}}[K]\) is \emph{not} canonically derived; it is introduced only to mark the distinction between kinematic closure and dynamical admissibility. The trefoil becomes a viable electron-like candidate only if
\begin{equation}
\delta E_{\mathrm{eff}}[T_{2,3}] = 0,
\qquad
\delta^2 E_{\mathrm{eff}}[T_{2,3}] > 0
\quad\text{(local minimum or long-lived metastable sector).}
\label{eq:trefoilstabilityTrefoil}
\end{equation}

\paragraph{Step 7b: The Topological Shielding Gate and Geometric Mass Limit.}
The energetic admissibility of the trefoil state is governed by the universal SST mass functional. To maintain a strictly hydrodynamic description, the invariant rest mass must decouple from electromagnetic coupling parameters. By evaluating the geometric limit, the inter-sector amplification factor is identified identically as the geometric impedance mismatch between the Compton boundary $\lambda_c$ and the vortex core radius $r_c$. The mass functional is thus defined entirely by spatial length ratios and topological invariants:
\begin{equation}
\label{eq:MassFunctional}
M(T_{2,3}) = \left( \frac{\lambda_c}{\pi r_c} \right)^{G(T)} \phi^{-g(T)} k(T)^{-3/2} n(T)^{-1/\phi} \left( \frac{\frac{1}{2}\rho_{\rm core} \lVert\mathbf{v}_{\!\boldsymbol{\circlearrowleft}}\rVert^2 \mathcal{V}_{\rm core}}{c^2} \right),
\end{equation}
where $\phi = (1+\sqrt{5})/2$ is the Golden Ratio controlling the discrete scale invariance of the medium, $g(T)$ is the Seifert genus, and $k(T)$ is the complexity index.

For the electron, the $2{:}3$ trefoil mode-lock triggers a binary topological shielding gate $G(T) = 1$, activating the macroscopic geometric impedance factor $\lambda_c / (\pi r_c)$. The fundamental distinction between continuous fluid dynamics and discrete particle physics in SST is thus reduced to this topological phase transition: without the exact $2{:}3$ rational lock, the configuration remains an unshielded fluid perturbation; with the lock, it inherits the discrete topological invariants $(k, n, g)$ that permanently fix its mass and prevent continuous decay.

\paragraph{Step 8: relation to the branch index \(\sigma\).}
The two-valued internal branch index of the main text should \emph{not} be identified with \(T_{2,3}\) versus \(T_{3,2}\), because those denote the same knot type under exchanged toroidal/poloidal naming. The more precise SST statement is
\begin{equation}
\boxedstatement{$\sigma=\pm 1$ labels the two internal handed circulation--torsion realizations of a fixed closure class, naturally associated with opposite chiral trefoils (mirror pair), not with $T_{2,3}$ versus $T_{3,2}$.}
\label{eq:boxedSigmaTrefoil}
\end{equation}
Equivalently, \(\sigma=+1\) and \(\sigma=-1\) are interpreted as the two opposite handed realizations of the trefoil closure sector relative to a chosen axis.

\paragraph{Step 9: connection back to the frozen interaction.}
This subsection does not replace the frozen interaction of the main text; it constrains the admissible circulation profile entering it. The circulation-based driver
\begin{equation}
\gamma(\mathbf r,t)=\frac{\Gamma_{\mathrm{loc}}(\mathbf r,t)}{\Gamma_0}
\label{eq:gammaFromTrefoil}
\end{equation}
is therefore not an arbitrary helical packet once trefoil closure is imposed. In the trefoil sector we denote the constrained field by
\begin{equation}
\gamma_{2{:}3}(\mathbf r,t)
\equiv
\gamma(\mathbf r,t)\Big|_{q/p=3/2\ \text{or}\ 2/3},
\label{eq:gamma23Def}
\end{equation}
with the branch label \(\sigma=\pm 1\) selecting the handed realization inside that trefoil-compatible sector. Thus the generic packet ansatz is restricted to the trefoil-compatible winding ratio and branch assignment. The frozen interaction remains
\begin{equation}
\delta H_{\mathrm{int}}^{(\sigma)}
=
-2\sigma E_R\eta_0\,\gamma
+
\frac{E_R\eta_0\,\sigma}{2}
\left\{\frac{\hat L_z}{\hbar},\gamma\right\},
\label{eq:HintTotalTrefoil}
\end{equation}
but the present subsection says that, in a trefoil-closure regime, the \(\gamma\) entering Eq.~\eqref{eq:HintTotalTrefoil} should be understood as the mode-locked circulation field \(\gamma_{2{:}3}\) rather than as a generic helical packet.

\begin{equation}
\boxedstatement{Trefoil closure does not replace the frozen interaction; it constrains the allowed circulation field $\gamma(\mathbf r,t)$ entering $\delta H_{\mathrm{int}}^{(\sigma)}$.}
\label{eq:boxedTrefoilConstrainsGamma}
\end{equation}

\paragraph{Step 10: working SST interpretation of spin and folding.}
At the present working level, the safer SST statement is
\begin{equation}
m_s=\pm\frac12
\quad\leadsto\quad
\sigma=\pm 1
\quad\leadsto\quad
\text{two internal handed circulation--torsion branches relative to a chosen axis,}
\label{eq:spinbranchworkingTrefoil}
\end{equation}
rather than a literal classical clockwise/counterclockwise mechanical spin. Then ``folding into a trefoil'' means: a closed circulation ring at the primitive turnover scale \(\Omega_0\) supports a helical torsion mode which phase-locks into the \(2{:}3\) closure ratio and is dynamically stable enough to persist.

\begin{equation}
\boxedstatement{The strongest defensible SST claim is not that the electron is already proven to be a trefoil, but that if the primitive internal rate is $\Omega_0 \approx \omega_c$, then trefoil realization requires a rational $2{:}3$ mode lock on a closed circulation ring.}
\label{eq:boxedStrongestClaim}
\end{equation}

\paragraph{Status.}
Eqs.~\eqref{eq:torusknotansatzTrefoil}--\eqref{eq:trefoilphaselockTrefoil} establish the kinematic trefoil criterion. Eqs.~\eqref{eq:EeffTopNormalized}--\eqref{eq:trefoilstabilityTrefoil} state the still-open dynamical admissibility condition without reintroducing unfrozen couplings. Eqs.~\eqref{eq:boxedSigmaTrefoil}, \eqref{eq:gamma23Def}, and~\eqref{eq:HintTotalTrefoil} connect the closure picture back to the branch index \(\sigma\) and to the frozen interaction already used in the v0.8 photoionization benchmark.

\paragraph{Analogy for a 10-year-old.}
A plain spinning ring just wiggles if you shake it. It becomes a trefoil only when the wiggle wraps around the ring in the right counting pattern --- like going around the donut \(2\) times while winding around its tube \(3\) times. The ring's own size creates a strict time delay, forcing the wave to either perfectly match its own tail and lock into this specific $2{:}3$ pattern, or quickly cancel itself out.

\section{Extensions and High-Energy Limits}
The present framework establishes the low-energy, incompressible atomic baseline. Several physical extensions naturally arise but are deferred to future work to preserve the strict phenomenological closure of the atomic Hamiltonian.

First, longitudinal compressional modes (acoustic phonons in the condensate) are excluded here, as their kinematic front speed $c_P \gg c$ decouples them from standard atomic spectra. Second, the explicit topological inheritance mechanisms during photon emission---whereby the helical torsion wave is modeled as an emitted skyrmion texture---remain to be formalized dynamically. Finally, vacuum magneto-optical effects, such as parity-odd Faraday rotation induced by the substrate's intrinsic chirality, are predicted as downstream consequences of the torsion field sector but are not strictly required to recover the Aufbau principle.

\section{Falsifiable Predictions and Symmetries}
The present model constitutes a viable physical theory only because its parameters are highly constrained. The core predictive commitments are:
\begin{enumerate}[label=(P\arabic*)]
    \item \emph{Hydrogenic recovery}: The Bohr radius and discrete energy spectrum emerge purely from thermodynamic Euler balance and NLS delay-locking.
    \item \emph{The Aufbau architecture}: The $2n^2$ shell capacity is derived fundamentally from the spatial acoustic limits of the medium coupled with the topological cavitation barrier.
    \item \emph{Branch symmetry}: In the absence of an external circulation driver ($\gamma=0$), the two chiral branches ($\sigma = \pm 1$) must remain strictly degenerate.
    \item \emph{Operator-level chirality}: The explicitly handed anticommutator interaction predicts an observable operator-level asymmetry in resonant photoionization cross-sections based on the incoming wave's topological winding.
    \item \emph{Fixed couplings}: The parameters $\Gamma_0$, $\gamma$, $\eta_0$, and $g_J$ are fixed analytically by the medium constants and allow no atom-by-atom phenomenological retuning.
\end{enumerate}

The strongest falsifier of this framework is the absence of chiral asymmetry in resonant interactions. If no branch-sensitive effects survive once the couplings are fixed, the branch interpretation reduces to a mere relabeling of ordinary spin. If branch-sensitive effects are observed, their predicted magnitude and symmetry scaling provide direct targets for attosecond photoemission experiments.

\section{Conclusion}
This framework demonstrates that quantum structural logic---including energy quantization, discrete orbital radii, and fermionic exchange---need not be postulated as irreducible laws of nature, but can emerge consistently from the continuous, nonlinear hydrodynamics of a locally structured vacuum.

By defining:
\begin{equation}
\boxed{
\text{propagating torsion}
+\text{ rotating-fluid sources}
+\text{ a branch-resolved Hamiltonian},
}
\end{equation}
we have established a sharply reduced, internally consistent atomic model. The derivation of the Pauli barrier from fluid cavitation (Appendix A) and the resulting Aufbau principle (Section 8) firmly anchor the physics of the periodic table in classical topology and continuum mechanics. This formulation provides a rigorous foundation that can be experimentally tested, extended to molecular bindings, or explicitly falsified term by term.

\paragraph{Analogy for a 10-year-old.} Imagine the atom as a globe with allowed storm patterns on it. A twisting wave packet can hit the globe and nudge one storm pattern into a higher allowed one. The two Swirl-Clock branches are like two opposite ways the storm can twist internally. The point of this paper is to say exactly which equations describe the wave, the circulation injection, and the storm pattern.

\appendix
\section{Topological Origin of Fermionic Exchange: The Hydrodynamic Pauli Barrier}
\label{app:PauliBarrier}

To extend the Swirl-String Theory (SST) atomic bridge beyond the hydrogenic baseline, the phenomenological Pauli exclusion principle must be derived from the foundational fluid mechanics of the vacuum. In the present framework, the electron is not a point particle subject to abstract anticommutation relations, but an extended $2{:}3$ mode-locked trefoil vortex loop carrying a quantized circulation $\Gamma_0$. The exclusion principle must therefore emerge as a strict hydrodynamic energy penalty prohibiting identical vortex states from occupying the same spatial domain.

\subsection{Biot-Savart Interaction Energy of Vortex Loops}
Consider two thin vortex loops, $C_1$ and $C_2$, embedded in an incompressible, inviscid fluid of effective density $\rho_{\!f}$. The total velocity field of the fluid is the linear superposition of the fields induced by each loop: $\mathbf{u}_{\rm total} = \mathbf{u}_1 + \mathbf{u}_2$. The total kinetic energy of the fluid is
\begin{equation}
\label{eq:Etotal_fluid}
E_{\rm total} = \frac{1}{2}\rho_{\!f} \int \lVert \mathbf{u}_1 + \mathbf{u}_2 \rVert^2 \, d^3r
= E_1 + E_2 + E_{\rm int},
\end{equation}
where $E_1$ and $E_2$ are the invariant self-energies (rest masses) of the isolated knots, and the interaction energy is
\begin{equation}
\label{eq:Eint_integral}
E_{\rm int} = \rho_{\!f} \int \mathbf{u}_1 \cdot \mathbf{u}_2 \, d^3r.
\end{equation}
By defining the vorticity $\boldsymbol{\omega}_i = \nabla \times \mathbf{u}_i$ and utilizing the vector potential $\mathbf{A}_i$ such that $\mathbf{u}_i = \nabla \times \mathbf{A}_i$, the incompressible constraint $\nabla \cdot \mathbf{u}_i = 0$ allows the vector potential to be written via the Biot-Savart kernel \cite{Batchelor1967}:
\begin{equation}
\label{eq:VectorPotential}
\mathbf{A}_i(\mathbf{r}) = \frac{1}{4\pi} \int \frac{\boldsymbol{\omega}_i(\mathbf{r}')}{\lVert \mathbf{r} - \mathbf{r}' \rVert} \, d^3r'.
\end{equation}
For thin vortex filaments where the vorticity is strictly confined to a core of radius $r_c$, the volume integral contracts to a line integral along the filament centerline. The vorticity element maps as $\boldsymbol{\omega}_i \, d^3r \to \Gamma_i \, d\mathbf{l}_i$. Substituting this back into Eq.~\eqref{eq:Eint_integral} via integration by parts yields the classical Kelvin interaction energy:
\begin{equation}
\label{eq:BiotSavartEnergy}
E_{\rm int} = \frac{\rho_{\!f} \Gamma_1 \Gamma_2}{4\pi} \oint_{C_1} \oint_{C_2} \frac{d\mathbf{l}_1 \cdot d\mathbf{l}_2}{\lVert \mathbf{r}_1 - \mathbf{r}_2 \rVert}.
\end{equation}

\subsection{The Topological Repulsion Penalty}
We now apply Eq.~\eqref{eq:BiotSavartEnergy} to two electrons in the same atomic orbital. Let both electrons correspond to the $2{:}3$ trefoil closure with the same canonical circulation magnitude $\Gamma_1 = \Gamma_2 = \Gamma_0 = 2\pi r_c \lVert\mathbf{v}_{\!\boldsymbol{\circlearrowleft}}\rVert$.

If the two electrons occupy the exact same spatial state and share the same internal Swirl-Clock branch ($\sigma_1 = \sigma_2$), their local tangent vectors $d\mathbf{l}_1$ and $d\mathbf{l}_2$ are completely parallel. Let $\mathbf{d}$ be the displacement vector between their geometric centers. As the two identical topological states are forced into the same location ($\mathbf{d} \to 0$), the denominator in Eq.~\eqref{eq:BiotSavartEnergy} approaches the fundamental cut-off scale $r_c$, and the dot product $d\mathbf{l}_1 \cdot d\mathbf{l}_2 = \lVert d\mathbf{l} \rVert^2 > 0$ is everywhere positive.

The interaction energy evaluates to a localized pseudo-potential:
\begin{equation}
\label{eq:VPauli}
V_{\rm Pauli}(\mathbf{d}) \approx \frac{\rho_{\!f} \Gamma_0^2}{4\pi} \oint \oint \frac{d\mathbf{l}_1 \cdot d\mathbf{l}_2}{\sqrt{\lVert \mathbf{r}_1(\theta_1) - \mathbf{r}_2(\theta_2) + \mathbf{d} \rVert^2 + r_c^2}}.
\end{equation}
When the loops are co-spatial ($\mathbf{d} \to 0$), constructive interference of their circulation fields creates an extreme dynamic pressure spike. The energetic penalty for full overlap evaluates approximately to the scale of the self-energy of the loop's macroscopic circumference, $L \approx 2\pi a_0^{\rm SST}$:
\begin{equation}
\label{eq:VPauli_Max}
V_{\rm Pauli}^{\rm max} \approx \frac{\rho_{\!f} \Gamma_0^2}{4\pi} \left( \frac{L}{r_c} \right) \mathcal{O}(1).
\end{equation}

Conversely, if the two electrons occupy the same spatial orbital but have \emph{opposite} Swirl-Clock branches ($\sigma_1 = -\sigma_2$), their circulation vectors are anti-parallel ($d\mathbf{l}_1 \cdot d\mathbf{l}_2 < 0$). This results in destructive interference of the local fluid velocity, leading to an attractive interaction ($E_{\rm int} < 0$), up to the point where the vortex cores would structurally recombine or annihilate.

Thus, the Pauli exclusion principle is not an independent postulate of SST; it is the direct classical consequence of the fluid kinetic energy integral. Identically handed, identically placed vortex knots repel with an insurmountable dynamic pressure barrier.

\subsection{Numerical Validation of the Exclusion Energy}
To rigorously validate this hydrodynamic exclusion mechanism, we compute the magnitude of the barrier using the canonical SST parameters:
\begin{align}
\rho_{\!f} &= 7.0 \times 10^{-7}\ \mathrm{kg\,m^{-3}}, \\
r_c &= 1.40897017 \times 10^{-15}\ \mathrm{m}, \\
\lVert\mathbf{v}_{\!\boldsymbol{\circlearrowleft}}\rVert &= 1.09384563 \times 10^6\ \mathrm{m\,s^{-1}}.
\end{align}
The circulation quantum is:
\begin{equation}
\Gamma_0 = 2\pi r_c \lVert\mathbf{v}_{\!\boldsymbol{\circlearrowleft}}\rVert \approx 9.6836 \times 10^{-9}\ \mathrm{m^2\,s^{-1}}.
\end{equation}
The dimensional energy prefix in Eq.~\eqref{eq:VPauli_Max} is $\rho_{\!f} \Gamma_0^2 \approx 6.564 \times 10^{-23}\ \mathrm{J \cdot m}$.

For two identical electron trefoils confined to the hydrogenic baseline orbital, the relevant macroscopic scale is the SST Bohr radius $a_0^{\rm SST} \approx 5.291 \times 10^{-11}\ \mathrm{m}$. Substituting these into the maximal overlap penalty yields:
\begin{equation}
\label{eq:VPauli_Num}
V_{\rm Pauli}^{\rm max} \approx \frac{6.564 \times 10^{-23}}{4\pi} \left( \frac{2\pi \times 5.291 \times 10^{-11}}{1.4089 \times 10^{-15}} \right) \approx 1.23 \times 10^{-18}\ \mathrm{J} \approx 7.6\ \mathrm{eV}.
\end{equation}

This result serves as a stringent falsification check. If the hydrodynamic Pauli penalty had evaluated to meV scales, it would be easily overcome by thermal noise, collapsing the atom. If it had evaluated to GeV scales, it would tear the atomic boundary apart. The derived topological repulsion barrier is explicitly on the order of $\sim 10\ \mathrm{eV}$---exactly the magnitude required to prevent identical orbital occupation while remaining safely bounded within the limits of atomic and chemical binding energies.

\subsection{The Cavitation Limit and the Absolute Exclusion Boundary}
While Eq.~\eqref{eq:VPauli_Num} establishes a large energetic barrier, the absolute nature of the Pauli exclusion principle in SST is enforced by the structural limits of the fluid itself. The incompressible medium possesses a strict upper bound on transmittable force, corresponding to the fluid's cavitation threshold, $F_{\rm swirl}^{\max} \approx 29.0535\ \mathrm{N}$.

The dynamic pressure drop induced by the constructive interference of two identically handed, co-spatial vortex cores scales as $\Delta P \sim -\frac{1}{2}\rho_{\!f} \lVert 2\mathbf{v}_{\!\boldsymbol{\circlearrowleft}} \rVert^2$. When the local interaction force exceeds $F_{\rm swirl}^{\max}$, the isotropic pressure of the condensate is overcome, leading to localized vacuum cavitation.

Therefore, fermionic exchange is geometrically forbidden not merely by an energy slope, but by a topological phase boundary: two identical electrons cannot occupy the same state because the resulting vorticity concentration would cavitate the ambient fluid substrate, destroying the local foliation structure required to sustain them.

\begin{thebibliography}{99}

\bibitem{Helmholtz1858}
H. von Helmholtz,
``Über Integrale der hydrodynamischen Gleichungen, welche den Wirbelbewegungen entsprechen,''
\textit{Journal für die reine und angewandte Mathematik} \textbf{55} (1858), 25--55.

\bibitem{Kelvin1869}
W. Thomson (Lord Kelvin),
``On Vortex Motion,''
\textit{Transactions of the Royal Society of Edinburgh} \textbf{25} (1869), 217--260.

\bibitem{Moffatt1969}
H. K. Moffatt,
``The degree of knottedness of tangled vortex lines,''
\textit{Journal of Fluid Mechanics} \textbf{35} (1969), 117--129.
DOI: 10.1017/S0022112069000991.

\bibitem{Batchelor1967}
G. K. Batchelor,
\textit{An Introduction to Fluid Dynamics},
Cambridge University Press, 1967.

\bibitem{Greenspan1968}
H. P. Greenspan,
\textit{The Theory of Rotating Fluids},
Cambridge University Press, 1968.

\bibitem{HehlObukhov2007}
F. W. Hehl and Y. N. Obukhov,
\textit{Elie Cartan's Torsion in Geometry and in Field Theory, an Essay},
Annales de la Fondation Louis de Broglie \textbf{32} (2007), 157--194.

\bibitem{Dirac1927}
P. A. M. Dirac,
``The Quantum Theory of the Emission and Absorption of Radiation,''
\textit{Proceedings of the Royal Society A} \textbf{114} (1927), 243--265.
DOI: 10.1098/rspa.1927.0039.

\bibitem{Dirac1928}
P. A. M. Dirac,
``The Quantum Theory of the Electron,''
\textit{Proceedings of the Royal Society A} \textbf{117} (1928), 610--624.
DOI: 10.1098/rspa.1928.0023.

\bibitem{Hartree1928}
D. R. Hartree,
``The Wave Mechanics of an Atom with a Non-Coulomb Central Field. Part I. Theory and Methods,''
\textit{Mathematical Proceedings of the Cambridge Philosophical Society} \textbf{24} (1928), 89--110.
DOI: 10.1017/S030500410001191X.

\bibitem{Pitaevskii2003}
L. Pitaevskii and S. Stringari,
\textit{Bose-Einstein Condensation},
Oxford University Press, 2003.

\bibitem{Erneux2009}
T. Erneux,
\textit{Applied Delay Differential Equations},
Springer, 2009.

\bibitem{Adler1946}
R. Adler,
``A Study of Locking Phenomena in Oscillators,''
\textit{Proceedings of the IRE} \textbf{34} (1946), 351--357.
DOI: 10.1109/JRPROC.1946.229930.

\bibitem{Kedia2013}
H. Kedia, I. Bialynicki-Birula, D. Peralta-Salas, and W. T. M. Irvine,
``Tying knots in light fields,''
\textit{Physical Review Letters} \textbf{111} (2013), 150404.
DOI: 10.1103/PhysRevLett.111.150404.

\bibitem{BerryDennis2001}
M. V. Berry and M. R. Dennis,
``Knotted and linked phase singularities in monochromatic waves,''
\textit{Proceedings of the Royal Society A} \textbf{457} (2001), 2251--2263.
DOI: 10.1098/rspa.2001.0826.

\end{thebibliography}

\end{document}